
\documentclass[11pt]{article}
\usepackage[a4paper,margin=0.95in]{geometry}
\usepackage{fontspec}
\setmainfont{Noto Serif}
\usepackage{microtype}
\usepackage{amsmath,amssymb,mathtools}
\usepackage{csquotes}
\usepackage{array}
\usepackage{hyperref}
\hypersetup{colorlinks=false,pdfborder={0 0 0}}
\setlength{\parindent}{1.15em}
\setlength{\parskip}{0.18em}
\allowdisplaybreaks
\sloppy
\newcommand{\scn}[1]{\textsc{#1}}
\newcommand{\arctg}{\operatorname{arctg}}
\begin{document}
\begin{center}
{\Large\bfseries ÜBER EINEN VON LEJEUNE DIRICHLET HERRÜHRENDEN BEWEIS AUS DER WAHRSCHEINLICHKEITSRECHNUNG\footnote{Das Folgende ist bis auf den zur Orientirung hinzugefügten ersten Absatz ein wörtlicher Abdruck aus dem Astronomischen Jahrbuch für 1834 S. 295--298; vgl. Encke, Gesammelte Abhandlungen Bd. II S. 49--53. F.}\par}
\vspace{0.3em}
{\large\bfseries Nach einer Mittheilung von J. F. Encke\par}
\vspace{0.8em}
\end{center}

Ordnet man die Beobachtungsfehler, ohne Rücksicht auf ihr Zeichen, nach ihrer absoluten Grösse, so wird bei $m$ Beobachtungen der, welcher zu dem Index $\frac12(m+1)$ gehört, bei ungeradem $m$, oder bei geradem $m$ das arithmetische Mittel zwischen den Fehlern mit dem Index $\frac12m$ und $\frac12m+1$ einen genäherten Werth für den wahrscheinlichen Fehler $r$ ergeben.

Wenn indessen schon bei den Potenzensummen die grössere Anzahl der Beobachtungsfehler die Genauigkeit in Bezug auf die wahrscheinlichen Grenzen so sehr wachsen lässt, so wird bei dieser Zählungsweise es um so mehr stattfinden müssen. Da \scn{Gauss} in der Zeitschrift für Astronomie Bd. I S. 195 die dazu nöthige Formel ohne Beweis angegeben, so wird um so mehr der folgende elegante Beweis, den ich der Mittheilung meines geehrten Collegen, Herrn Prof. \scn{Dirichlet}, verdanke, hier von Werth sein, als der Satz selbst anderswo noch nicht bewiesen ist.

Man suche die Wahrscheinlichkeit, dass bei $(2n+1)$ Beobachtungen die Vertheilung der Fehler so sei, dass ein Fehler liege zwischen $t$ und $t+dt$, $n$ Fehler zwischen $0$ und $t$, und $n$ Fehler zwischen $t+dt$ und $\infty$. Die Wahrscheinlichkeit, dass ein Fehler kleiner als $t$ sei, sei wiederum ganz allgemein
\[
 \int_0^t\psi(\Delta)d\Delta=u.
\]
Die Wahrscheinlichkeit eines Fehlers $>t+dt$ wird dann werden
\[
 1-\psi(t)dt-\int_0^t\psi(\Delta)d\Delta=1-u-\psi(t)dt,
\]
da die Wahrscheinlichkeit eines Fehlers zwischen $t$ und $t+dt$ ist $=\psi(t)dt$.

Hiernach ist die zusammengesetzte Wahrscheinlichkeit einer Anordnung der Fehler, wenn $n$ Fehler $<t$, ein Fehler zwischen $t$ und $t+dt$, und $n$ Fehler $>t+dt$,
\[
 =u^n(1-u)^n\psi(t)dt,
\]
wenn man die Glieder zweiter Ordnung vernachlässigt, da das Resultat von der ersten Ordnung ist. Solcher Fälle oder Anordnungen können aber so viele vorkommen als Versetzungen von $2n+1$ Elementen möglich sind, wenn unter ihnen $n$ gleiche Elemente einer Art (deren Wahrscheinlichkeit $=u$) und $n$ gleiche Elemente einer andern Art (deren Wahrscheinlichkeit $=(1-u)$) vorkommen. Folglich ist die Wahrscheinlichkeit aller möglichen Anordnungen dieser Art
\[
 =\frac{1\cdot2\cdot3\cdots(2n+1)}{(1\cdot2\cdot3\cdots n)^2}u^n(1-u)^n\psi(t)dt=U.
\]
Denkt man sich die Grösse $dt$ des Intervalls zwischen $t$ und $t+dt$ constant, so giebt es einen Werth von $t$, für welchen $U$ ein Maximum ist. Die sich durch Differentiation zur Bestimmung desselben ergebende Gleichung ist:
\[
 \frac{n\psi(t)}{u}-\frac{n\psi(t)}{1-u}+\psi'(t)=0,
\]
wo $\psi'(t)$ die nämliche Bedeutung, wie oben $\varphi'(\Delta)$ hat. Es ist nämlich $du$, oder das Increment von $\int_0^t\psi(\Delta)d\Delta$ in Bezug auf eine unendlich kleine Aenderung der Grenze $t$, gleich $\psi(t)dt$. Man kann der letzten Gleichung die Form geben
\[
 \frac1u-\frac1{1-u}+\frac{\psi'(t)}{n\psi(t)}=0.
\]
Das letzte Glied wird um so kleiner werden, je grösser $n$ ist, oder je mehr Beobachtungen gegeben sind. Bei einer hinlänglich grossen Anzahl wird man es vernachlässigen können. Oder der Werth von $t$, für welchen das Maximum stattfindet, nähert sich bei wachsendem $n$ immer mehr dem Werthe, welcher aus der Gleichung folgt:
\[
 \frac1u-\frac1{1-u}=0,
\]
d. h.
\[
 u=\int_0^t\psi(\Delta)d\Delta=\frac12,
\]
oder nach der oben gegebenen Definition dem Werthe $r$.

Nimmt man das Integral von $U$ zwischen bestimmten Grenzen, so erhält man daraus die Wahrscheinlichkeit, dass der in der Mitte liegende Fehler in diesen Grenzen enthalten ist. Diese wird also für die Grenzen $r-\delta$ und $r+\delta$
\[
 =\frac{1\cdot2\cdot3\cdots(2n+1)}{(1\cdot2\cdot3\cdots n)^2}\int_{r-\delta}^{r+\delta}u^n(1-u)^n\psi(t)dt;
\]
oder weil $\psi(t)dt=du$, wenn wegen der Grenzen in Bezug auf $t$ gesetzt wird
\[
 \int_0^{r-\delta}\psi(t)dt=u',\qquad \int_0^{r+\delta}\psi(t)dt=u'',
\]
so wird die Wahrscheinlichkeit, dass der mittelste Werth zwischen $r-\delta$ und $r+\delta$ liegt,
\[
 =\frac{1\cdot2\cdot3\cdots(2n+1)}{(1\cdot2\cdot3\cdots n)^2}\int_{u'}^{u''}u^n(1-u)^n\,du.
\]
Je grösser die Anzahl der Beobachtungen ist, desto enger werden die Grenzen, zwischen welchen $t$ mit gleicher Wahrscheinlichkeit liegen wird. Sind deshalb die Beobachtungen zahlreich genug, so wird, wenn man $u'$ und $u''$ nach dem Taylor'schen Satze entwickelt, es erlaubt sein, nur die erste Potenz von $\delta$ zu berücksichtigen. Dadurch wird:
\[
 u'=\int_0^r\psi(t)dt-\delta\psi(r)\cdots=\frac12-\delta\psi(r)
\]
und ebenso
\[
 u''=\frac12+\delta\psi(r).
\]
Sowohl diese Form, als auch die Verbindung von $u$ und $1-u$ in dem Integral, zeigt an, dass man eine noch bequemere Form erhalten wird, wenn man für $u$ eine andere Variable einführt; am besten durch die Gleichung
\[
 u=\frac12+\frac{s}{2\sqrt n}=\frac12\left(1+\frac{s}{\sqrt n}\right),
\]
folglich
\[
 1-u=\frac12-\frac{s}{2\sqrt n}=\frac12\left(1-\frac{s}{\sqrt n}\right),
\]
wobei die Grenzen in Bezug auf $s$ gefunden werden durch
\[
 \delta\psi(r)=\frac{s}{2\sqrt n}.
\]
Hiernach wird das Integral
\[
 \frac{1\cdot2\cdot3\cdots(2n+1)}{(1\cdot2\cdot3\cdots n)^2}\cdot\frac{1}{2^{2n+1}\sqrt n}\int_{-2\delta\sqrt n\psi(r)}^{+2\delta\sqrt n\psi(r)}\left(1-\frac{s^2}{n}\right)^n ds,
\]
oder weil $s$ im Differential nur gerade Potenzen enthält,
\[
 \frac{1\cdot2\cdot3\cdots(2n+1)}{(1\cdot2\cdot3\cdots n)^2}\cdot\frac{1}{2^{2n}\sqrt n}\int_0^{2\delta\sqrt n\psi(r)}\left(1-\frac{s^2}{n}\right)^n ds.
\]
Sei nun $\delta\sqrt n$ eine endliche Grösse $=\varrho$, also die Grenze $\delta$ in eben dem Maasse abnehmend wie $\sqrt n$ zunimmt, so bleibt $s$ innerhalb der angenommenen Grenze endlich, wie sehr auch $n$ zunimmt. Bei einem grossen $n$ wird man aber auch nach der Entwickelung der Logarithmen in \scn{Euler}'s \emph{Introductio} setzen können:
\[
 \left(1-\frac{s^2}{n}\right)^n=e^{-s^2}
\]
und
\[
 \frac{1\cdot2\cdot3\cdots2n}{(1\cdot2\cdot3\cdots n)^2}=\frac{2^{2n}}{\sqrt{n\pi}},
\]
nach \scn{Euler} \emph{Calc. Diff.} P. II. Cap. VI. \S 160--162, als dem Grenzwerthe, welchem es sich beständig mit wachsendem $n$ nähert, so dass der Ausdruck wird
\[
 \frac{2n+1}{n\sqrt\pi}\int_0^{2\delta\sqrt n\psi(r)}e^{-s^2}\,ds,
\]
wofür man unbedenklich schreiben kann
\[
 \frac2{\sqrt\pi}\int_0^{2\delta\sqrt n\psi(r)}e^{-s^2}\,ds
\]
als den Ausdruck für die Wahrscheinlichkeit, dass bei zahlreichen Beobachtungen der mittelste Fehler $r$, wenn alle der Grösse nach geordnet sind, liegt zwischen $r-\delta$ und $r+\delta$.

Diese Wahrscheinlichkeit wird folglich $\frac12$, oder es sind die wahrscheinlichen Grenzen gegeben durch
\[
 2\delta\sqrt n\,\psi(r)=\varrho,
\]
d. h.
\[
 \delta=\frac{\varrho}{2\sqrt n}\cdot\frac1{\psi(r)}.
\]
Für das oben angenommene Gesetz der Fehler
\[
 \psi(\Delta)=2\varphi(\Delta)=\frac{2h}{\sqrt\pi}e^{-hh\Delta\Delta}
\]
werden also die wahrscheinlichen Grenzen von $r$
\[
 r\pm\frac{\varrho e^{hhr r}\sqrt\pi}{4\sqrt n\,h},
\]
oder wenn man statt $2n+1$ die Anzahl der Beobachtungen $m$ nennt, und die Gleichung
\[
 hr=\varrho
\]
benutzt:
\[
 r\left\{1\pm\frac{e^{\varrho\varrho}\sqrt\pi}{\sqrt{8m}}\right\}.
\]
Der numerische Werth von $e^{\varrho\varrho}$ ist $1{,}2554176$, womit der Ausdruck wird:
\[
 r\left\{1\pm\frac{0{,}786716}{\sqrt m}\right\}.
\]

\end{document}
